\section{Roadmap Update after Experiment 12A}
\label{sec:roadmap-post-12a}

Experiment~12A clears the nonlinear-learning viability gate but does \emph{not} clear the practical-certificate scaling gate. This distinction is important.

The supported conclusions are now:
\begin{enumerate}
  \item \textbf{Sparse temporal event communication survives a nonlinear convex learning objective.} With an appropriately rescaled trigger, scheduled coordinate events approach the centralized logistic-regression test loss using only a few thousand logical payload bits.
  \item \textbf{The descent-aware oracle remains valuable outside quadratics.} The exact global descent oracle again reaches the centralized reference with very low communication and substantially better transient optimization than the heuristic jump schedule.
  \item \textbf{Valid curvature certificates remain safe.} Full, block, and diagonal-plus-residual summaries produced zero measured harmful accepted events in the strong non-IID safety audit.
  \item \textbf{The present practical summaries are too expensive or too conservative.} Full and large-block certificates spend most of their communication on dense gradient calibration and curvature metadata. Small summaries reject too many useful events.
  \item \textbf{The diagonal-only control identifies the remaining opportunity.} It is numerically strong and much cheaper, but is not a valid certificate because cross-coordinate drift is ignored. Its harmful-event rate shows that the missing coupling information matters.
\end{enumerate}

Therefore the project should \textbf{not} proceed directly to the small-neural-network Experiment~13. Doing so would reproduce the same unresolved certificate-information bottleneck at much larger dimension and make the failure harder to diagnose.

\subsection{Recommended next step: Experiment 12B}

The next study should be \textbf{Experiment~12B: selective block calibration and compressed coupling certificates}. Its purpose is to close the gap between the safe but expensive full/block certificates and the cheap but unsafe diagonal-only rule.

Two communication costs must be attacked simultaneously:
\begin{enumerate}
  \item the one-time curvature/coupling summary;
  \item the recurring dense gradient calibration.
\end{enumerate}
Experiment~12A shows that reducing only the first is insufficient because recurring calibration already dominates the total payload for the useful safe operating points.

\subsection{12B.1: selective coordinate/block calibration}

The most direct extension is to let different coordinates or blocks maintain different calibration ages. For block $B$, the server stores the calibrated gradient components $g_B^{\rm cal}$ and the model state $w_B^{\rm cal,model}$ at which those gradient components were refreshed. The full current model is already known at the server, so storing the calibration model incurs no communication.

When an event in block $B$ cannot be certified, clients should refresh only the gradient components required by that block rather than sending all $d$ components. If $b=\abs{B}$, one such refresh costs approximately
\begin{equation}
  B_{\rm block\,cal}=N b\times32
\end{equation}
bits instead of $Nd\times32$. For the current $N=10$, $d=20$ problem, a five-coordinate block costs $1600$ payload bits rather than $6400$.

The certificate remains valid because each coordinate $j$ can maintain its own calibration snapshot $w_j^{\rm cal}$ and use
\begin{equation}
  \abs{\nabla_jF(w)-g_j^{\rm cal}}
  \leq
  \sum_k\bar M_{jk}\abs{w_k-w_{j,k}^{\rm cal}}.
\end{equation}
Thus asynchronous calibration ages are not a theoretical obstacle.

\subsection{12B.2: coupling-aware blocks}

The contiguous blocks used in Experiment~12A are only a neutral baseline. A better block partition should group strongly coupled coordinates, for example using the weighted graph
\begin{equation}
  A_{jk}=\bar M_{jk}
\end{equation}
or a normalized variant. Blocks can then be chosen to minimize the residual cross-block coupling that makes the uncertainty envelope conservative.

The key comparison should be
\begin{equation}
  \text{contiguous blocks}
  \quad\text{vs}\quad
  \text{coupling-clustered blocks}.
\end{equation}
This is the practical counterpart of the basis-alignment lesson from Experiments~10D and~11B.

\subsection{12B.3: compressed residual coupling}

If block summaries remain too conservative, a low-rank-plus-residual representation should be tested. Let $A_i\geq0$ be a compressed approximation of $M_i$. A valid bound can always be restored by defining the nonnegative residual
\begin{equation}
  R_i=\max(M_i-A_i,0)
\end{equation}
componentwise. Then
\begin{equation}
  M_i\leq A_i+R_i.
\end{equation}
The structured part $A_i$ can be represented by low-rank factors or block factors, while only a residual row-sum bound needs to be retained for the unresolved coupling. This creates an explicit communication--certificate-tightness trade-off.

\subsection{Experiment 12B success criterion}

Experiment~12B should be considered successful only if a \emph{valid} selective/block certificate produces a useful Pareto point relative to both:
\begin{enumerate}
  \item the very cheap scheduled-event baseline;
  \item EF-TopK.
\end{enumerate}
In particular, the method should preserve the zero-harmful-event behavior of the valid Experiment~12A certificates while reducing calibration/curvature payload enough that improved transient or final performance is not purchased by an order-of-magnitude communication increase.

If selective block calibration succeeds, the project can proceed to Experiment~13 with layer/block certificates. If it fails, the safest scientific conclusion is that sparse temporal event communication itself is useful, but rigorous global-descent certification requires too much side information to scale in its current form.

The updated progression is therefore
\begin{center}
\resizebox{0.96\linewidth}{!}{$\displaystyle \boxed{12\mathrm{B}\;\text{selective/block certificate}\;\rightarrow\;13\;\text{small neural network only if the Pareto gate passes}.}$}
\end{center}
